
\subsubsection{Problem Setup}

Consider a classification task where inputs x have class label y and group membership g. Groups arise from the interaction of core features (predictive of y) and spurious features (correlated with y in training but not causally related). Under ERM training, models exploit spurious correlations because they provide simpler, more consistent signals.

\subsubsection{Loss-Based Detection}

For each sample i, we track loss $L_i$(t) across training epochs t. We use loss at an early detection epoch T\_early = 5 (5\% of 100 epochs). Candidate minority samples are identified as those with loss above the k-th percentile:

minority\_candidate\_i = 1[$L_i$(T\_early) > percentile\_k(L(T\_early))]

At 5\% minority rate, k = 95 selects the top 5\% of high-loss samples.

\subsubsection{Linear Probe Analysis}

To verify the simplicity bias mechanism, linear probes are trained on frozen ResNet-18 features at multiple epochs. Two separate logistic regression probes are trained:

\begin{itemize}
\item \textbf{Spurious probe}: Predicts background (water/land)
\item \textbf{Core probe}: Predicts bird type (waterbird/landbird)
\end{itemize}

If simplicity bias operates, spurious probe accuracy should exceed core probe accuracy at early epochs.

\subsubsection{Implementation Details}

\begin{itemize}
\item \textbf{Model}: ResNet-18 with ImageNet pretrained weights (IMAGENET1K\_V1)
\item \textbf{Training}: 100 epochs, SGD (lr=0.001, momentum=0.9, weight decay=1e-4), batch size 64
\item \textbf{Checkpoints}: Epochs 5, 20, 50, 81, 100
\item \textbf{Seed}: 42
\end{itemize}

